Much of the solving algorithms detailed later in this manual utilizes finite volume method to solve for the desired states within a droplet. Since the droplets can deviate from perfect sphericity and each turn into
an ellipsoid, the use of an ellipsoidal coordinate system is appropriate. There are many ways to formulate such system; the one used here resembles spherical coordinates, but with a different sketch factor in different
semi-axial directions.

\begin{equation}
    \mathbf{r}(r, \mu, \varphi) =
    \begin{bmatrix}
        x' \\ y' \\ z'
    \end{bmatrix}
    = \begin{bmatrix}
        a r \mu\\
        b r \sqrt{1-\mu^2} \cos\varphi\\
        c r \sqrt{1-\mu^2} \sin\varphi
    \end{bmatrix},
\end{equation}
where $r \in [0, 1]$, $\mu \in [-1, 1]$, and $\varphi \in (-\pi, \pi]$ are the coordinate variables. The axes $x'$, $y'$, and $z'$ form a rotated Cartesian grid where $x'$ points to the
semi-major direction, $y'$ to the semi-middle direction, and $z'$ to the semi-minor direction. As such, the polar cosine $\mu$ is always measured from the semi-major axis.

The reverse transformations are
\begin{subequations}
    \begin{align}
        r &= \sqrt{\left( \frac{x'}{a} \right)^2 + \left( \frac{y'}{b} \right)^2 + \left( \frac{z'}{c} \right)^2 } \\
        \mu &= \frac{x'}{ar} \\
        \varphi &= \arctantwo\left( \frac{z'}{c}, \frac{y'}{b} \right)
    \end{align}
\end{subequations}

The basis vectors are defined as followed
\begin{subequations}
    \begin{align}
        \mathbf{e}_r &\equiv \frac{\partial\mathbf{r}}{\partial r} = \frac{\mathbf{r}}{r} \\
        \mathbf{e}_\mu &\equiv \frac{\partial\mathbf{r}}{\partial \mu} =
            \begin{bmatrix}
                ar \\ \displaystyle -br \frac{\mu}{\sqrt{1-\mu^2}} \cos\varphi \\ \displaystyle -cr \frac{\mu}{\sqrt{1-\mu^2}} \sin\varphi
            \end{bmatrix} \\
        \mathbf{e}_\varphi &\equiv \frac{\partial\mathbf{r}}{\partial \varphi} =
            \begin{bmatrix}
                0 \\ -br \sqrt{1-\mu^2} \sin\varphi \\ cr \sqrt{1-\mu^2} \cos\varphi
            \end{bmatrix}.
    \end{align}
\end{subequations}

The normalized counterparts are denoted as $\mathbf{\hat{r}}$, $\bm{\hat{\upmu}}$, and $\bm{\hat{\upvarphi}}$.